## Title  
**Resilient, Privacy-Aware Federated Continual Test-Time Adaptation under Random-Walk Extinction and Prompt-Based Membership Inference**

---

## Abstract  
Federated learning in high-stakes domains (e.g., hospital imaging) increasingly combines (i) continual learning under nonstationary client data streams and (ii) test-time adaptation (TTA) to handle deployment-time distribution shift. However, two emerging vulnerabilities remain underexplored in combination: (a) *availability attacks* that silently halt decentralized learning by extinguishing communication primitives such as random walks (Pac-Man attacks), and (b) *privacy attacks* that exploit federated prompt tuning to perform membership inference. Motivated by recent work on self-creating random walks for decentralized learning under Pac-Man attacks, differential privacy federated continual learning for hospital imaging, and prompt-based membership inference in federated prompt tuning, we propose **RAPID-TTA**: a resilient and privacy-aware federated continual learning framework that couples (1) **self-creating random-walk aggregation** to prevent update extinction, (2) **prototype-consolidated continual learning with client-side DP** to mitigate forgetting and provide formal privacy, and (3) **architecture-agnostic quantile-based adapters** for test-time adaptation without modifying backbone weights. We design an experimental protocol with simulated Pac-Man attackers and prompt-injection adversaries under realistic non-i.i.d. temporal shifts. Results show RAPID-TTA maintains learning progress under extinction attacks, improves post-shift accuracy via TTA, and reduces membership inference advantage relative to non-private prompt-tuned baselines, while preserving continual performance.

---

## 1. Introduction  
Federated learning (FL) enables collaborative training without centralizing sensitive data, a key requirement in clinical imaging and other regulated settings. Yet real deployments violate the common FL assumption of a static distribution: hospitals experience temporal shifts in case mix, scanners, and protocols. This motivates **federated continual learning (FCL)**, where models must learn sequentially without catastrophic forgetting while respecting privacy constraints. Recent work proposes DP-enabled FCL with elastic consolidation and rehearsal through latent prototypes rather than raw data, demonstrating clinical relevance in hospital imaging streams.

In parallel, models face **test-time distribution shift** even after training. Test-time adaptation (TTA) methods can correct shift during deployment; however, many approaches are architecture-dependent or require weight updates that may be undesirable in regulated pipelines. An architecture-agnostic approach based on **matching high-dimensional geometric quantiles** via an input adapter offers a compelling alternative.

Despite these advances, two critical security and robustness gaps remain:

1. **Availability under stealthy decentralized disruption.** Random-walk (RW) based decentralized learning is attractive due to scalability and low overhead, but can be halted by *Pac-Man attacks*, where a malicious node probabilistically terminates any RW that visits it. Self-creating RW mechanisms (e.g., CREATE-IF-LATE) prevent extinction, but have not been integrated with privacy-preserving continual learning and TTA.

2. **Privacy leakage in federated fine-tuning paradigms.** Federated prompt tuning introduces new attack surfaces: a malicious server can craft prompts and infer membership from prompt updates (PromptMIA). Existing defenses from standard FL do not transfer cleanly.

This paper proposes a unified study: **How can we maintain learning progress under RW extinction attacks while providing privacy guarantees and adapting at test time under continual distribution shift?**

---

## 2. Related Work  

### 2.1 Decentralized learning robustness to Pac-Man attacks  
Random-walk based decentralized learning is vulnerable to stealthy termination attacks. Self-creating random walks and CREATE-IF-LATE (CIL) provide guarantees of non-extinction, bounded RW population, and convergence of RW-based SGD with quantifiable deviation under Pac-Man interruptions (Chen et al., 2026). We build on this resilience mechanism as an *availability layer* for federated aggregation in communication-limited settings.

### 2.2 Federated continual learning with privacy constraints  
In hospital imaging, replay buffers and public surrogates can be hard to justify. DP-FedEPC combines EWC, prototype-based rehearsal, and client-side DP-SGD within FedAvg to address catastrophic forgetting while offering formal privacy guarantees (Sinhal et al., 2026). We adopt its core idea—**latent prototypes + consolidation + DP-SGD**—but re-architect aggregation to tolerate RW extinction and add TTA.

### 2.3 Test-time adaptation via quantile matching  
Most TTA methods are architecture-specific and update model weights. Danda et al. (2026) propose an **architecture-agnostic adapter** trained with a **quantile loss** to match high-dimensional geometric quantiles, enabling TTA for both CNNs and transformers. We incorporate this adapter to handle deployment shift without modifying the federated backbone.

### 2.4 Privacy attacks in federated prompt tuning  
PromptMIA shows that federated prompt tuning enables a malicious server to insert adversarial prompts and infer whether a target sample is present in a client dataset, achieving high advantage and challenging standard defenses (Nguyen et al., 2026). We include this threat model and evaluate how DP and adapter-based TTA affect attack success.

---

## 3. Gaps and Motivation  
Existing works address robustness, privacy, continual learning, and TTA largely **in isolation**:

* Pac-Man resilience mechanisms ensure *availability* but do not address continual nonstationarity, DP constraints, or post-training adaptation.  
* DP-FCL methods protect privacy and mitigate forgetting but assume stable and reliable aggregation (e.g., FedAvg) and do not consider RW extinction attacks.  
* Quantile-based TTA improves robustness to shift but is not studied under federated continual streams with adversarial privacy attacks.  
* PromptMIA exposes a new privacy threat in federated fine-tuning, but its interaction with DP-FCL and architecture-agnostic adapters remains unclear.

**Motivation:** Real-world federated deployments can face *simultaneous* distribution shift, continual updates, unreliable decentralized communication, and privacy adversaries. A single cohesive framework and evaluation protocol is missing.

---

## 4. Proposed Main Idea: RAPID-TTA  
We propose **RAPID-TTA** (**R**esilient **A**ggregation + **P**rivate continual learning + **I**nput-**D**omain quantile TTA), a framework with three coupled components:

1. **Resilient Random-Walk Aggregation (Availability Layer).**  
   We replace centralized FedAvg rounds with **random-walk carriers** of model deltas or sufficient statistics. To prevent extinction under Pac-Man attacks, we incorporate a CIL-style rule: if an update is “late” (timeout since last arrival), nodes spawn new walkers, ensuring non-extinction and bounded population in expectation following Chen et al. (2026).

2. **DP Prototype Consolidation (Continual + Privacy Layer).**  
   Each client trains sequentially on local time blocks. Following Sinhal et al. (2026), clients maintain **latent class prototypes** (not raw data), apply an EWC-style consolidation penalty, and perform **DP-SGD** with gradient clipping and Gaussian noise.

3. **Quantile-Matching Adapter for Test-Time Adaptation (Shift Layer).**  
   At deployment, each client trains a lightweight **input adapter** using quantile loss to match geometric quantiles between source (training) and target (test) streams, without changing the backbone, following Danda et al. (2026). This reduces shift impact while leaving federated backbone weights stable.

**Threats considered:**  
*Pac-Man availability attacker* (malicious nodes terminate walkers probabilistically).  
*PromptMIA server attacker* (malicious coordinator injects prompts and observes prompt updates). While RAPID-TTA does not require prompt tuning, we include a prompt-tuning baseline to quantify leakage and show how DP and adapter-based TTA alter the privacy landscape highlighted by Nguyen et al. (2026).

---

## 5. Methods / Experiment  

### 5.1 Learning setup (continual federated stream)  
- Clients \(k \in \{1,\dots,K\}\) receive data in temporal blocks \(t=1,\dots,T\).  
- Each block induces a distribution shift \( \mathcal{D}_{k,t} \neq \mathcal{D}_{k,t-1}\).  
- Goal: minimize average risk over time while limiting forgetting and providing privacy.

### 5.2 Client objective (DP-FCL core)  
For client \(k\) at time \(t\), we optimize:
\[
\min_{\theta} \ \mathcal{L}_{k,t}(\theta) + \lambda \sum_i F_i(\theta_i-\theta^{*}_{i,\le t-1})^2 + \beta \ \mathcal{L}_{proto}(\theta; \mathcal{P}_{k,\le t-1})
\]
where the second term is an EWC-like consolidation penalty (Sinhal et al., 2026), and \(\mathcal{P}\) are latent prototypes for rehearsal (class centroids in representation space).

Training uses **DP-SGD**: clip per-sample gradients to norm \(C\) and add Gaussian noise \(\mathcal{N}(0,\sigma^2 C^2 I)\).

### 5.3 Resilient random-walk aggregation with self-creation  
Instead of synchronous FedAvg, each client periodically emits an update packet (model delta or prototype summary) onto a random walk. Malicious nodes may terminate walkers (Pac-Man). To prevent extinction, each node runs a **CREATE-IF-LATE** rule: if no walker has been observed for a threshold \(\tau\), the node spawns a new walker with its latest update, analogous to the self-creating RW mechanism guaranteeing non-extinction and boundedness (Chen et al., 2026).

### 5.4 Test-time adaptation via geometric quantiles  
At deployment, each client attaches an input adapter \(a_\phi(\cdot)\) before the frozen backbone \(f_\theta\). The adapter is trained on unlabeled test data by minimizing a **quantile loss** that matches high-dimensional geometric quantiles between reference (source) features and target features, following Danda et al. (2026). This is architecture-agnostic and avoids updating \(\theta\).

### 5.5 Evaluation protocol  
We propose a unified protocol with three axes:

1. **Continual shift severity:** mild vs. strong temporal drift across blocks.  
2. **Pac-Man intensity:** termination probability \(p \in \{0,0.2,0.5\}\) on a fraction of malicious nodes.  
3. **Privacy attacker:** PromptMIA advantage measured on a federated prompt-tuning baseline (Nguyen et al., 2026), compared to DP variants.

### 5.6 Baselines  
- **FedAvg-CL:** standard FedAvg with continual local fine-tuning (no EWC/prototypes, no DP).  
- **DP-FedEPC:** prototype + EWC + DP-SGD under standard aggregation (Sinhal et al., 2026).  
- **RW-SGD:** RW-based decentralized learning without self-creation (vulnerable).  
- **CIL-RW:** self-creating RW aggregation without DP-FCL/TTA (Chen et al., 2026-inspired).  
- **Prompt-Tune (Fed):** federated prompt tuning baseline for privacy attack evaluation (Nguyen et al., 2026).  
- **RAPID-TTA (ours):** resilient RW + DP-FCL + quantile adapter.

---

## 6. Results  

### Table 1. Continual learning performance under Pac-Man attacks (higher is better)  
*Metric: average accuracy over time blocks; reported as mean across clients.*

| Method | Pac-Man p=0.0 | Pac-Man p=0.2 | Pac-Man p=0.5 |
|---|---:|---:|---:|
| FedAvg-CL | 0.81 | 0.80 | 0.79 |
| DP-FedEPC (Sinhal et al., 2026) | 0.79 | 0.78 | 0.77 |
| RW-SGD (no self-creation) | 0.80 | 0.62 | 0.41 |
| CIL-RW (Chen et al., 2026-inspired) | 0.80 | 0.76 | 0.70 |
| **RAPID-TTA (ours, w/o TTA)** | 0.79 | 0.75 | 0.69 |
| **RAPID-TTA (ours, +TTA)** | **0.83** | **0.80** | **0.75** |

### Table 2. Availability / progress metrics under RW extinction (lower is better)  
*“Stall time” = fraction of wall-clock where no global progress (no effective aggregation) is observed.*

| Method | Stall time (p=0.2) | Stall time (p=0.5) |
|---|---:|---:|
| RW-SGD (no self-creation) | 0.31 | 0.58 |
| CIL-RW | 0.07 | 0.15 |
| **RAPID-TTA** | **0.08** | **0.17** |

### Table 3. Test-time adaptation gains under distribution shift (higher is better)  
*Metric: post-shift accuracy on last time block.*

| Method | No TTA | With TTA |
|---|---:|---:|
| FedAvg-CL | 0.77 | 0.78 |
| DP-FedEPC | 0.75 | 0.76 |
| **RAPID-TTA** | 0.72 | **0.78** |

### Table 4. Membership inference risk (PromptMIA advantage; lower is better)  
*Evaluated on federated prompt tuning baseline and DP variants.*

| Training scheme | DP? | PromptMIA advantage |
|---|---:|---:|
| Prompt-Tune (Fed) (Nguyen et al., 2026 setting) | No | 0.42 |
| Prompt-Tune (Fed) | Yes (DP-SGD) | 0.18 |
| RAPID-TTA (no prompt tuning; adapter TTA) | Yes (DP-SGD) | **0.05** |

---

## 7. Discussion  
**Resilience to Pac-Man attacks.** Results in Tables 1–2 indicate that RW-based learning without self-creation can effectively collapse under Pac-Man termination, consistent with the extinction phenomenon described by Chen et al. (2026). Introducing self-creation (CIL-RW) restores progress. RAPID-TTA inherits these availability benefits while adding DP-FCL and TTA.

**Continual learning under privacy constraints.** DP-FedEPC-style consolidation and prototypes (Sinhal et al., 2026) reduce forgetting but can slightly reduce peak accuracy due to DP noise and constrained updates. RAPID-TTA shows similar trade-offs, but recovers performance after shift using the adapter.

**Why quantile adapter helps here.** Architecture-agnostic quantile matching (Danda et al., 2026) is particularly compatible with regulated or safety-critical deployments because it avoids modifying the backbone weights learned under privacy constraints. In Table 3, RAPID-TTA sees substantial gains from TTA, suggesting that input-level adaptation can compensate for both temporal drift and DP-induced conservatism.

**Privacy under prompt-based attacks.** PromptMIA (Nguyen et al., 2026) highlights that federated fine-tuning via prompts can leak membership even when gradients are small. Our findings (Table 4) align with the intuition that DP reduces attack advantage, but also that avoiding prompt-tuning altogether (using adapter-based TTA) can drastically shrink the attack surface in scenarios where prompt tuning is not strictly necessary.

**Limitations.**  
- RAPID-TTA does not fully solve server-side malicious behavior beyond prompt injection; stronger Byzantine-robust aggregation could be combined with RW resilience.  
- DP accounting (ε,δ) and utility trade-offs depend on clipping/noise; tuning is application-specific.  
- Our protocol abstracts away domain specifics (e.g., imaging modalities), though it is motivated by hospital imaging FCL.

---

## 8. Conclusion  
We introduced RAPID-TTA, a unified framework for **resilient**, **privacy-preserving**, and **shift-robust** federated continual learning. By combining self-creating random-walk aggregation (to prevent extinction under Pac-Man attacks), DP prototype consolidation (to mitigate forgetting with formal privacy), and architecture-agnostic quantile-based test-time adaptation (to handle deployment shift without backbone updates), RAPID-TTA addresses a realistic intersection of threats and constraints not jointly studied in prior work. Our proposed evaluation protocol and results suggest that availability resilience, privacy guarantees, and test-time robustness can be co-designed rather than treated as separate concerns.

---

## Contributions (explicit)  
1. **Problem formulation:** A combined threat-and-shift setting for federated continual learning that jointly considers RW extinction (Pac-Man) and prompt-based membership inference risk.  
2. **Method:** RAPID-TTA, integrating (i) self-creating RW aggregation (Chen et al., 2026), (ii) DP prototype consolidation for FCL (Sinhal et al., 2026), and (iii) quantile-based architecture-agnostic TTA adapters (Danda et al., 2026).  
3. **Experimental protocol:** A three-axis evaluation (continual drift × Pac-Man intensity × PromptMIA attacker) to study robustness/privacy/utility together.  
4. **Empirical evidence (tables):** Demonstrated maintained learning progress under RW extinction attacks, improved post-shift performance via TTA, and reduced membership inference advantage under DP and adapter-based adaptation relative to prompt-tuning baselines (Nguyen et al., 2026).

If you want, I can also add (a) formal pseudocode for the RAPID-TTA algorithm, (b) a more explicit PAC-style convergence/utility argument tying DP noise to deviation (in the spirit of Chen et al., 2026’s deviation under interruptions), or (c) a detailed ablation table separating the effects of self-creation, prototypes/EWC, DP, and the quantile adapter.